%\medskip
%
%\begin{tabularx}{0.8\linewidth}{|X|}\hline
%\textbf{Pour chaque question, recopier sur la copie son numéro et la réponse
%correspondante.}\\
%\textbf{Pour cette partie, aucune justification n’est demandée.}\\
%\textbf{Pour les questions a choix multiple, une seule réponse est exacte.}\\ \hline
%\end{tabularx}
%\end{center}

\medskip

\textbf{Question 1}

%Quelle est la mesure, en degrés, d’un angle droit ?
90~degrés

\medskip

\textbf{Question 2}

%Voici une série de quatre notes : 8~;~10~;~11~;~11. 

%Quelle est la moyenne de cette série ?
%
%\begin{center}
%\begin{tabularx}{\linewidth}{*{4}{X}}
%\textbf{A.~~}9,5 &\textbf{B.~~}10 &\textbf{C.~~}10,5 &\textbf{D.~~}11
%\end{tabularx}
%\end{center}
%\medskip
La moyenne est égale à $\dfrac{8 + 10 + 11 + 11}{4} = \dfrac{40}{4} = 10$ : réponse B.
\medskip

\textbf{Question 3}

%Dans un collège de 800 élèves, 25\,\% des élèves portent des lunettes.
%
%Combien d’élèves portent des lunettes ?

25\,\% des 800 élèves représente $\dfrac{25}{100} \times 800 = \dfrac{25 \times 800}{100} = 25 \times 8 = 200$~(élèves).

\emph{Remarque } : $\dfrac{25}{100} = \dfrac{25 \times 1}{4 \times 25} = \dfrac14$ : c'est le quart, soit la moitié de la moitié : on part de 800 et on divise deux fois par 2 : $800 \longmapsto 400 \longmapsto 200$.

\medskip

\textbf{Question 4}

%Le graphique ci-dessous donne l’évolution de la température (en degrés Celsius) en fonction de l’horaire (en heures).
%
%\begin{center}
%\psset{xunit=0.5cm,yunit=0.125cm}
%\begin{pspicture}(-2,-5)(18,35)
%\multido{\n=0+2}{10}{\psline[linewidth=0.2pt](\n,0)(\n,35)}
%\multido{\n=0+5}{8}{\psline[linewidth=0.2pt](0,\n)(18,\n)}
%\psaxes[linewidth=1.25pt,Dx=2,Dy=5,labelFontSize=\scriptstyle](0,0)(0,0)(18,35)
%\psdots[dotscale=1.5](8,15)(12,27)(16,30)
%\psline(8,15)(12,27)(16,30)
%\uput[d](9,-3){Horaire (en heures)}
%\rput{90}(-2,17.5){Température (en \degres C)}
%\end{pspicture}
%\end{center}

%Entre 8~h et 16~h, de combien de degrés la température a-t-elle augmenté ?
La température est passée de $15^\circ$ à $30^\circ$ soit une élévation de température de $15^\circ$.

\medskip

\textbf{Question 5}

%Une voiture roule à 90~km/h. Combien de temps met-elle pour parcourir 45~km ?~
%\begin{center}
%\begin{tabularx}{\linewidth}{*{4}{X}}
%\textbf{A.~~}15 min &\textbf{B.~~}30 min &\textbf{C.~~}45~min &\textbf{D.~~}1~h
%\end{tabularx}
%\end{center}
La voiture parcourt 90 km en 1 heure, donc 45 km en 0,5~h = 30~min
\medskip

%\begin{minipage}{0.65\linewidth}
\textbf{Question 6}

%Donner le périmètre du losange ABCD représenté ci- contre.
%\end{minipage}\hfill
%\begin{minipage}{0.32\linewidth}
%\psset{unit=1cm}
%\begin{pspicture}(-2.5,-2.2)(2.5,2.2)
%\uput[l](-2.3,0){A} \uput[u](0,2){D} \uput[r](2.3,0){C} \uput[d](0,-2){B} \uput[ul](-1.15,1){3 cm} 
%\pspolygon(-2.3,0)(0,2)(2.3,0)(0,-2)%ABCD
%\end{pspicture}
%\end{minipage}
Le losange a ses quatre côtés de même longueur ; son périmètre est égal à $4 \times 3 = 12$~(cm).

\medskip

\textbf{Question 7} (1 point)

%Pour résoudre l’équation $4x - 3 = 20$, on effectue le calcul :
%
%\begin{center}
%\begin{tabularx}{\linewidth}{*{4}{X}}
%\textbf{A.~~}$x = \dfrac{20}{4} + 3$ &\textbf{B.~~}$x = (20 - 4) + 3$ &\textbf{C.~~}$x = 20 \times 4 + 3$ &\textbf{D.~~}$x = \dfrac{20 + 3}{4}$
%\end{tabularx}
%\end{center}
Pour résoudre l'équation $4x - 3 = 20$, on ajoute 3 à chaque membre :

\[4x = 20 + 3\]

4 est un facteur : on le neutralise pae son inverse $\dfrac14$ :

\[\dfrac14 \times 4x = \dfrac14\left(20 + 3\right) \:\text{ou } \: x  = \dfrac{(20 + 3)}{4}.\]

Réponse D

\medskip

\textbf{Question 8} (1 point)

%\begin{minipage}{0.55\linewidth}
%Sur la figure ci-contre, les droites (DE) et (AC) sont parallèles.
%
%Écrire une égalité de rapports permettant de déterminer la longueur AB.
%
%\end{minipage}\hfill
%\begin{minipage}{0.4\linewidth}
%\psset{unit=1cm}
%\begin{pspicture}(5.4,4)
%\pspolygon(0.3,0.1)(4.7,0.6)(0.4,3.4)%ACB
%\psline(0.35,1.2)(3.37,1.5)%DE
%\uput[dl](0.3,0.1){A}  \uput[r](4.7,0.6){C}\uput[ul](0.4,3.4){B}
%\uput[l](0.35,1.2){D}\uput[ur](3.37,1.5){E}
%\uput[d](2.5,0.35){6}\uput[d](1.86,1.35){4}\uput[ur](1.88,2.45){4,5}
%\uput[l](0.375,2.3){3}
%\end{pspicture}
%\end{minipage}
%\medskip
$\bullet~$Les points B, D et A sont alignés dans cet ordre ;

$\bullet~$Les points B, E et C sont alignés dans cet ordre ;

$\bullet~$les droites (DE) et (AC) sont parallèles.

On a donc une configuration de Thalès et les mesures des côtés des triangles BAC et BDE sont proportionnelles : 

en particulier $\dfrac{\text{BA}}{\text{BD}} = \dfrac{\text{AC}}{\text{DE}}$, soit 
$\dfrac{\text{BA}}{3} = \dfrac{6}{4}$ ; on en déduit en multipliant chaque membre par 3 : BA $= 3 \times \dfrac64 = \dfrac{3 \times 2 \times 3}{2 \times 2} = \dfrac92 = 4,5$.

\medskip

\textbf{Question 9} (1 point)

On obtient successivement : $1 \longmapsto 1 \times 8 = 8 \longmapsto 8 + 10 = 18 \longmapsto \dfrac{18}{2} = 9$.

%\medskip
%
